\chapter{Case Studies}\label{ch:cases}

Three systems in this repository show the reduction at work on protocols nobody wrote
to be small: a multiplayer game server, a file store replicated across islands that may
be cut off from each other for a day, and LeanFM's own MCP protocol. The measurements
not generated for this book come from the experiments in
\texttt{reports/ctl-translation/faithful}, which carry their scripts.

\section{GoServer: sessions that fork and join}

GoServer's protocol is the composition Chapter~\ref{ch:where} described: sign-in mints a
person's session; listing the room's games is polled beside the lobby; a game runs as a
session of its own; a seat belongs to both a person and a game and runs four kinds of
request side by side.

\paragraph{Size comes from products.} As one model, GoServer has 39{,}497 checker nodes
and 1{,}604 typestates, with lookahead $2$ (Table~\ref{tab:corpus}). Almost all of it is
interleaving: the polling branch beside the lobby, and the seat's four child sessions.
Checked one session at a time (Section~\ref{prop:sessions}), every GoServer protocol
has lookahead $0$ except the seat, which keeps its own \texttt{parallel}: 941 nodes, 412
typestates, lookahead $1$.

\paragraph{A real log.} Two people signing in, one creating a game, both taking seats
and moving (Chapter~\ref{ch:where}): 24 events, no session ids, seven sessions derived,
the game a child of its creator's session, every seat and move a join of a person and
the game. The log conforms.

\paragraph{Faithfulness.} GoServer's typestates are not deterministic, so LeanFM decides
faithfulness by refinement: the protocol against itself in 1.42\,s through the server,
and a copy whose refused sign-in falls through to the room caught in the same time, with
the path to the difference (Chapter~\ref{ch:deciding}).

\section{IslandUploads: where the hidden state is}

IslandUploads stores files on islands that keep accepting writes while partitioned and
replicate once healed. Its ten request--response protocols (upload, update, feed,
download, the replication pulls, derived jobs, a forwarding door and a websocket) are
small: nine have lookahead $0$ with three to nine typestates each, and the websocket,
whose two directions are independent, has 38 typestates and lookahead $1$. Mutation
over their graphs caught every one of 792 changed copies and flagged none of 400
bisimilar ones.

Its architecture models are different. They describe the whole system: two islands, a
partition that starts and heals, commits, pulls across the link, feed order, import
lanes, conflict policy and retries. Here the lookahead reaches $5$ and $6$, and the
rules, written as trees, total about $6\times 10^9$ nodes for the unordered
replication model, against $14{,}757$ distinct subformulas. The reason is the one
Chapter~\ref{ch:constructs} found in bounded loops: the messages do not say what they
carry. After \texttt{PullAtoB}, which writes island B now holds depends on everything
before it, and a formula can only recover that by looking ahead.

The experiments measured each remedy on the cursor-based replication design
(typestates as LeanFM computes them, rules as Section~\ref{def:rules} writes them):

\begin{center}
\begin{tabular}{@{}lrr@{}}
\toprule
What each message carries & Lookahead & Rules as trees \\
\midrule
labels only, as written & 5 & $7.1\times 10^{8}$ nodes \\
the timeout also names its direction & 5 & $1.4\times 10^{9}$ nodes \\
each pull also names the records it delivered & 2 & $1.7\times 10^{6}$ nodes \\
a prev pointer, so each typestate is read off the chain & 0 & $59{,}502$ nodes \\
\bottomrule
\end{tabular}
\end{center}

\noindent Payload fields shorten the lookahead but cannot remove it: one message on a
shared timeline cannot say what the other island has done. A prev pointer removes it:
every rule becomes the one-step shape, total size is linear in the model's steps, and
faithfulness is decided by the walk (Section~\ref{thm:walk}), which on GoServer's graph
with prev labels took 17\,ms in the experiments. That is the design lesson of this book
in one table: \emph{a protocol whose messages carry their place in it is cheap to check
exactly.}

\section{LeanFM's own protocol}

LeanFM is developed with LeanFM, and its MCP protocol (\texttt{LeanFM.leanfm}: an agent, its host, the view, the
server and the disk; initialising, listing tools, calling a session tool and drawing its
answer, paging through project files, and closing) is the first thing
its own tools are tested on. It has 109 checker nodes, 35 typestates and lookahead $1$,
and the faithfulness check agrees with itself:

\ToolOutput{generated/ctl-leanfm-self.tex}

Its tests check the whole chain of this book on it. Every one of its 36 faithful rules
is proved by \texttt{session\_ctl} over the protocol. Over more than a hundred broken
copies of its graph (each step dropped, and steps sent to other states, relabelled and
added) the walk says ``faithful'' exactly when the CTL checker proves every rule, and the
walk and the refinement agree on every copy. That is the faithfulness result (Section~\ref{thm:faithful}),
the walk (Section~\ref{thm:walk}) and the cost of refinement (Section~\ref{prop:refine}) checked against each other, on
the protocol of the tool that checks them.
